\chapter{Zahlendarstellung}

\section{Natürliche Zahlen}

\subsection{Ursprünge der Zahlendarstellung}
Die \textit{natürlichen Zahlen} werden so genannt, da sie auf \enquote{natürliche Weise} beim Zählen verwendet werden.
Auf dem Gebiet der heutigen Demokratischen Republik Kongo wurde im 20. Jahrhundert ein Knochen, der sogenannte \textit{Ishango-Knochen}, gefunden.
Dieser Knochen wird auf die Zeit vor etwa 20'000 Jahren vor unserer Zeitrechnung datiert.
In den Knochen wurden deutlich erkennbar von Menschen einige Kerben eingeritzt.
Die Bedeutung dieser Kerben ist unklar, es wird jedoch angenommen, dass sie eine bestimmte Anzahl festhalten sollten.

Stellen wir uns vor, dass die Kerben die Anzahl der gefangenen Fische darstellen.
Gefangene Fische durch Kerben in einem Knochen zu repräsentieren, verlangt bereits einen recht hohen Grad an
Abstraktion --- schliesslich haben gefangene Fische und Kerben in einem Knochen auf den ersten Blick keinen offensichtlichen Zusammenhang.
Weniger abstrakt wäre es, eine Darstellung zu wählen, welche direkt an das zu zählende Objekt gebunden ist.
In diesem Fall würde man also keine Kerben ritzen, sondern die entsprechende Anzahl Fische zeichnen / skizzieren.
Eine Anzahl von drei Fischen könnte also durch drei Fisch-Symbole $\twemoji{tropical fish}\twemoji{tropical fish}\twemoji{tropical fish}$
festgehalten werden.

Anstelle von Kerben in einem Knochen oder Fisch-Symbole als Markierung könnte eine bestimmte Anzahl Fische auch durch
die entsprechende Anzahl von anderen Markierungen wie Striche ($|$) oder Kieselsteine (\textbullet) repräsentiert werden:
\begin{equation*}
	\text{drei Kerben} \quad\leftrightarrow\quad \twemoji{tropical fish}\twemoji{tropical fish}\twemoji{tropical fish} \quad\leftrightarrow\quad |||
	\quad\leftrightarrow\quad \text{\textbullet\textbullet\textbullet}
\end{equation*}
Die Darstellung durch skizzierte Fische ist am wenigsten abstrakt, da die dabei verwendete Markierung den zu zählenden Objekten (in diesem Fall also den Fischen) direkt nachempfunden ist.
Zahlendarstellungen, welche nur ein Symbol / eine Markierung (Kerbe, \twemoji{tropical fish}, $|$, \textbullet) verwenden, werden
\tib{unäre Zahlendarstellungen}\index{Zahlendarstellung!unäre} genannt.

\begin{myremark}\phantom{1}
  \begin{enumerate}[(a)]
    \item Der Vorteil von unären Darstellungen (gegenüber \enquote{komplexeren} Darstellungen) ist ihre einfache Verständlichkeit und
offensichtliche Interpretation.
    \item Unäre Darstellungen sind nur für kleine natürliche Zahlen übersichtlich, da ihre Darstellungslängen linear mit der Grösse der zu
darstellenden Zahl wachsen (doppelt so grosse Zahl $\rightarrow$ doppelt so lange Darstellung).
    \item Unäre Darstellungen eignen sich nicht besonders gut zum Rechnen.
  \end{enumerate} 
\end{myremark}

\subsection{Basisgrössen}
Anstelle von nur einem einzigen Symbol, wie bei unären Darstellungen, haben bereits Kulturen der Antike unterschiedliche
\tib{Basisgrössen}\index{Basisgrössen} verwendet. Die Römer verwendeten die sieben verschiedenen Basisgrössen
\begin{center}
  $I$ (Eins), $V$ (Fünf), $X$ (Zehn), $L$ (Fünfzig), $C$ (Hundert), $D$ (Fünfhundert) und $M$ (Tausend). 
\end{center}
\begin{myexample}\label{myexample:roman}
  Die Zahl \textit{fünfhundertvierundfünfzig} könnte in den römischen Basisgrössen zum Beispiel als
  \begin{equation*}
    C + C + C + C + C + L + I + I + I + I
  \end{equation*}
  dargestellt werden. 
\end{myexample}
Wir werden uns jede Basisgrösse als einen Typ von Münze von einem bestimmten Wert vorstellen. Die darzustellende Zahl
denken wir uns dann als einen bestimmten Geldbetrag, den wir mithilfe dieser gegebenen Basisgrössen (Münztypen) auszahlen wollen. In
dieser Vorstellung hatten die Römer also ein Währungssystem von sieben verschiedene Münztypen. Der Geldbetrag \textit{einundzwanzig} könnte
dann beispielsweise durch $X + X + I$ (zwei Zehnermünzen und eine Einermünze) in antiker römischer Währung ausbezahlt werden. Die alten Ägypter hingegen verwendeten damals bereits Symbole (Zeichnungen) für die Basisgrössen $1, 10, 100, 1000, 10000,
100000$ und $1000000$.

Ein Geldbetrag kann sowohl im römischen als auch ägyptischen System typischerweise auf mehrere unterschiedliche Arten
ausbezahlt werden (zum Beispiel, in dem nur Einer verwendet werden). Sie kennen diese Situation vom Bezahlen mit
Bargeld --- so kann der Betrag von hundert Schweizer Franken auf viele verschiedene Arten mit Schweizer Bargeld bezahlt werden
(zwei Fünfzigernoten, fünf Zwanzigernoten, eine Fünfzigernote $+$ zwei Zwanzigernoten $+$ zwei Fünfliber). Natürlich ist
es bei Zahlendarstellungen vorteilhaft, wenn sie eindeutig sind. Angenommen wir haben einen Betrag $x$ gegeben, welcher durch
unterschiedliche Ansammlungen von Münzen ausbezahlt werden kann. Dann betrachten wir diejenige Ansammlung von Münzen als
die \tib{Zahlendarstellung}\index{Zahlendarstellung} von $x$, welche am wenigsten Münzstücke benötigt.

\begin{myexample}\label{myexample:not_unique}
  Der Betrag \textit{sieben} kann im römischen System auf genau zwei verschiedene Arten (abgesehen von der Reihenfolge) ausbezahlt werden:
  \begin{itemize}
    \item $I + I + I + I + I + I + I$ (benötigt sieben Münzstücke),
    \item $V + I + I$ (benötigt drei Münzstücke).
  \end{itemize}
  Die Darstellung $V + I + I$ verwendet dabei am wenigsten Münzstücke und wird somit als die Zahlendarstellung von
  \textit{sieben} verstanden.
\end{myexample}
Nun gibt es aber denkbare Münzsysteme, in denen nicht jeder Betrag auf eindeutige Weise mit möglichst wenigen
Münzstücken ausbezahlt werden kann. Dies werden Sie in \cref{myexercise:1001} untersuchen.

\begin{myexercise}\label{myexercise:1001}
  Erweitern Sie das römische Münzsystem um zwei weitere Münztypen Ihrer Wahl und finden Sie einen Betrag, der in Ihrem
  neuen System mehrere unterschiedliche Darstellungen mit derselben minimalen Anzahl von Münzen erlaubt.
  \begin{myanswer}
    Wir erweitern das römische Münzsystem um die beiden Münztypen \coin{II} und \coin{III}. Dann besitzt der Betrag
    \textit{vier} die beiden Darstellungen $I$ $+$ \coin{III} und \coin{II} $+$ \coin{II}, welche beide eine (minimale) Anzahl von zwei Münzstücken verwenden.
  \end{myanswer}
\end{myexercise}

\section{Stellenwertdarstellung}
Das uns vermutlich am besten vertraute Zahlensystem ist das \tib{Dezimalsystem}\index{Stellenwertsystem!Dezimalsystem}
(Zehnersystem). Es verwendet die \textbf{zehn} verschiedenen Ziffern
\begin{align*}
  \lrc{0,1,2,3,4,5,6,7,8,9}
\end{align*}
und die (natürlichen) Potenzen $10^0, 10^1, 10^2, \ldots$ der Basis $10$ als \tib{Stellenwerte}\index{Stellenwert}.
Bereits in der Primarschule haben Sie gelernt, dass beispielsweise das Symbol $\green{4}\orange{6}\blue{0}\red{5}$ definiert ist als die folgende
Summe:
\begin{center}
  \red{fünf} Einer ($\red{5}\cdot 10^0$) $+$ \blue{null} Zehner ($\blue{0}\cdot 10^1$) $+$ \orange{sechs} Hunderter
  ($\orange{6}\cdot 10^2$) + \green{vier} Tausender
  ($\green{4}\cdot 10^3$). 
\end{center}
Eine Ziffer gibt dabei jeweils an, wie viele Male die entsprechende Potenz der Basis $10$ verwendet werden soll. Das
Dezimalsystem hat also (die unendlich vielen verschiedenen) Basisgrössen
\begin{align*}
  10^0=1,\quad 10^1=10,\quad 10^2=100,\quad 10^3=1000,\ldots
\end{align*}

\begin{mydefinition}\label{mydefinition:Stellenwertdarstellung}
  Seien $a_0, a_1, \ldots, a_n$ natürliche Zahlen. Dann ist der Ausdruck
  \begin{align*}
    a_na_{n-1}\ldots a_1a_0
  \end{align*}
  definiert durch die Summe
  \begin{align*}
    a_n\cdot 10^n + a_{n-1}\cdot 10^{n-1}+\ldots + a_1\cdot 10^1 + a_0\cdot 10^0.
  \end{align*}
\end{mydefinition}

\begin{myexample}
  Der Ausdruck $4205$ entspricht der Zahl
  \begin{align*}
    \textcolor{red}{4205} := \textcolor{red}{\overbrace{4}^{=a_3}}\cdot \underbrace{10^3}_{=1000} + \textcolor{red}{\overbrace{2}^{=a_2}}\cdot \underbrace{10^2}_{=100} + \textcolor{red}{\overbrace{0}^{=a_1}}\cdot \underbrace{10^1}_{=10} + \textcolor{red}{\overbrace{5}^{=a_0}}\cdot \underbrace{10^0}_{=1}.
\end{align*}
\end{myexample}

Anstelle der Potenzen der Basis $10$ kann man auch andere natürliche Zahlen als Basis verwenden\dots
\begin{myexercise}
  Begründen Sie, warum die Wahlen der natürlichen Zahlen $0$ oder $1$ nicht sinnvoll sind.
\end{myexercise}

Das \tib{Binärsystem}\index{Stellenwertsystem!Binärsystem} (Zweiersystem) verwendet die Ziffern $\lrc{0,1}$ und die
(natürlichen) Potenzen der $2$ als Basisgrössen.
